% This is samplepaper.tex, a sample chapter demonstrating the
% LLNCS macro package for Springer Computer Science proceedings;
% Version 2.25 of 2026/09/03
%
\documentclass[runningheads]{llncs}
%
\usepackage[T1]{fontenc}
% T1 fonts will be used to generate the final print and online PDFs,
% so please use T1 fonts in your manuscript whenever possible.
% Other font encondings may result in incorrect characters.
%

\usepackage{newtxtext}       % 
\usepackage[varvw]{newtxmath}       % selects Times Roman as basic font


\usepackage{graphicx}
% Used for displaying a sample figure. If possible, figure files should
% be included in EPS format.
%
% If you use the hyperref package, please uncomment the following two lines
% to display URLs in blue roman font according to Springer's eBook style:
%\usepackage{color}
%\renewcommand\UrlFont{\color{blue}\rmfamily}
%\urlstyle{rm}
%
\usepackage{comment}
\begin{document}
\renewcommand{\tablename}{Tabla}
%
\title{Evaluación de cercos virtuales dinámicos para la detección de merodeo peatonal y trayectorias recurrentes mediante visión artificial en videovigilancia inteligente.}
%
\titlerunning{Videovigilancia inteligente para detección de merodeo}
% If the paper title is too long for the running head, you can set
% an abbreviated paper title here
%
\author{Denilson Celi\inst{1}\orcidID{0009-0009-4480-2770} \and
Raúl Hinojoza\inst{1}\orcidID{0009-0009-9635-3292}}
%
\authorrunning{D. Celi et al.}
% First names are abbreviated in the running head.
% If there are more than two authors, 'et al.' is used.
%
\institute{Universidad de las Fuerzas Armadas ESPE, Sangolquí 171103, Ecuador \\
\email{\{dsceli, rfhinojoza\}@espe.edu.ec}}
%
\maketitle              % typeset the header of the contribution
%
%\begin{abstract}
%.
 
%\keywords{\and }
 
%\end{abstract}
%
%
%
\section{Introducción}

Los Sistemas Inteligentes de Vigilancia (ISS) han experimentado una adopción global masiva producto del rápido avance en infraestructura tecnológica y la creciente demanda de seguridad en espacios públicos \cite{choudhry2023,wahyono2023}.


Mehmood \cite{mehmood2021} explica cómo, tradicionalmente, la vigilancia era ejercida por operadores humanos encargados de observar y analizar las transmisiones de video en pantallas, pero dado que los eventos anómalos eran poco frecuentes por naturaleza, esta labor se tornaba lenta, laboriosa e ineficiente, sobre todo en periodos prolongados de tiempo. Es en este contexto que surge la necesidad de contar con sistemas para la detección automática de anomalías. Los ISS nacen con el fin de suplir esta necesidad, introduciendo técnicas de visión por computador para el análisis de las imágenes generadas por las cámaras y con ello identificar sucesos anómalos, y todo este proceso con una intervención humana mínima \cite{wahyono2023}.

Entre los comportamientos que estos sistemas buscan detectar se encuentra el merodeo, el cual, aplicado al escenario peatonal, se define como un patrón observable de permanencia prolongada o desplazamiento continuo de un individuo en un área específica, sin una trayectoria o propósito evidentes, siendo normalmente asociada con actividades predelictivas como el robo  \cite{mehmood2021,wahyono2023,waddenkery2024}. En la identificación de esta anomalía resulta igual de fundamental el análisis de las trayectorias recurrentes, es decir, los patrones latentes que se obtienen de los puntos iniciales y finales de las rutas donde los individuos aparecen, pasan o se detienen reiteradamente en el espacio de la imagen \cite{jeoung2026}. Al capturar y clasificar estas trayectorias repetitivas, el sistema tiene la capacidad de determinar áreas de transición funcional y analizar el desplazamiento del individuo en su dominio espacial (posición del sujeto) y en su dominio temporal (tiempo que dura el evento) para discernir si su conducta constituye una amenaza \cite{jeoung2026,wahyono2023}. 

Para tratar los comportamientos anómalos mencionados previamente, las primeras versiones de analítica usaron algoritmos de visión por computadora que se basan en características creadas manualmente: el empleo de histogramas de gradientes orientados y descriptores para examinar las trayectorias en términos espaciales y temporales; la separación del fondo estático de los objetos en movimiento mediante modelos gaussianos mixtos \cite{mehmood2021,wahyono2023,waddenkery2024}; y la categorización de los eventos a través de modelos logísticos y lineales. La inflexibilidad del sistema ante factores externos, como las oclusiones, los cambios de la iluminación o los ángulos de cámara, y la necesidad de numerosas calibraciones complejas eran las principales desventajas de este enfoque original \cite{choudhry2023}.

En las fases iniciales de aplicación del aprendizaje profundo en analítica de video, propiciadas por una mayor capacidad computacional y la disponibilidad de conjuntos de datos más extensos, se implementaron arquitecturas convolucionales complementadas con algoritmos para rastrear vehículos y personas cuadro a cuadro \cite{vinod2024,nayak2019}. Para analizar secuencias temporales en video, se incorporaron redes convolucionales 3D, consistentes en dos ramas que combinan datos espaciales e información temporal, además de integrar redes recurrentes para captar las dependencias a través del tiempo \cite{mehmood2021,choudhry2023}. Dado que el refinamiento de estos sistemas resultaba complejo y costoso por el procesamiento y etiquetado masivo de horas de video, tecnologías como los autoencoders y las redes generativas antagónicas ofrecieron una alternativa viable, modelando patrones de conducta normal para emitir alertas únicamente cuando ocurre alguna desviación o anomalía \cite{choudhry2023,cho2025}. Sin embargo, aunque estas arquitecturas reducen la dependencia de datos completamente anotados, continúan operando sobre la totalidad de la escena visual. Este procesamiento integral obliga al algoritmo a evaluar ruido de fondo y áreas carentes de interés, sosteniendo una carga computacional ineficiente reflejada principalmente en el consumo de CPU, GPU, memoria y energía \cite{vinod2024}.

Estas limitaciones en los sistemas de videovigilancia inteligente evidencian que el procesamiento integral de la escena de video satura los recursos del sistema al obligarlo a evaluar ruido de fondo continuo y áreas irrelevantes. Con el propósito de reducir esta demanda computacional, se plantea la delimitación algorítmica del análisis visual mediante cercos virtuales dinámicos, los cuales permiten al usuario definir y reajustar manualmente los límites del área de interés según los requerimientos de vigilancia de la escena. Por consiguiente, el objetivo de esta investigación es evaluar el impacto de los cercos virtuales dinámicos sobre la precisión de clasificación y el costo computacional frente al procesamiento de escena completa en modelos de detección orientados al merodeo peatonal, para determinar su viabilidad algorítmica en tiempo real.



%
%
%
\section{Trabajos relacionados}

La literatura reciente sobre videovigilancia inteligente, seleccionada bajo la metodología PRISMA 2020 de Page et al. \cite{page2021} a partir de Web of Science, IEEE Xplore y Scopus como se resume en la Figura 1, se centra en la identificación de merodeo o trayectorias peatonales repetitivas y anómalas \cite{cho2025}, la localización y el seguimiento continuo de peatones en cámaras fijas mediante visión artificial o aprendizaje profundo, y la delimitación geométrica o semántica del espacio con líneas de cruce, regiones de interés o zonas restringidas configurables.

\begin{figure}[htbp]
\centering
\includegraphics[width=0.6\textwidth]{img/PRISMA.png}
\caption{Diagrama de flujo de selección bibliográfica según la metodología PRISMA 2020.}\label{fig1}
\end{figure}

Dentro de los estudios orientados a la detección de merodeo mediante seguimiento peatonal continuo, una primera línea metodológica combina modelos de detección de objetos y algoritmos de rastreo con reglas basadas en umbrales de tiempo o trayectoria definidos de antemano. En esta línea, Shanmughapriya et al. \cite{m2022} integraron YOLOv3, DeepSORT y clasificadores Haar Cascade en un entorno residencial para declarar el merodeo cuando un mismo peatón rastreado supera los 60 segundos de permanencia, mientras que Zhang et al. \cite{zhang2023} ampliaron dicho criterio en el sistema JYLD mediante YOLOA-S y JYTracker al combinar el tiempo de estancia con umbrales de distancia acumulada en píxeles y giros superiores a 90° y 160°. Pese a que el segundo trabajo incorpora restricciones geométricas para describir mejor el desplazamiento, ambos estudios coinciden en no reportar métricas de exactitud ni tasas de falsas alarmas sobre conjuntos de datos anotados, limitando sus pruebas a medir el tiempo de respuesta durante fotogramas consecutivos en vivo \cite{m2022} o a una inspección visual cualitativa \cite{zhang2023}. Esto dificulta conocer su precisión real al evaluar toda la escena sin delimitar zonas de interés.

A diferencia de las pruebas meramente cualitativas, otros estudios incorporan evaluaciones estadísticas del merodeo sobre conjuntos de datos anotados. Li y Huang \cite{li2024} midieron esta conducta al integrar YOLOX con un rastreador de siluetas de marcha y un criterio de desvío angular, reportando una exactitud de 87.4\,\% en el conjunto IITB-Corridor frente a métodos basados en esqueletos como MPED-RNN, que alcanzó 80.2\,\% en el mismo conjunto y presentó fallos ante desplazamientos en sentido opuesto \cite{jiang2022,li2024}. Sin embargo, en esta comparación los umbrales angulares se calibraron directamente sobre las secuencias de prueba, lo que introduce un sesgo de ajuste manual ligado a cada escena. En contraste, Waddenkery y Soma \cite{waddenkery2024} emplearon validación cruzada de 10 pliegues sobre el conjunto UCF-Crime alcanzando 99.76\,\% de exactitud, mientras que Mehmood \cite{mehmood2021} evaluó las métricas de forma desglosada por clase mediante redes convolucionales 3D logrando 98.5\,\% de exactitud. Estos dos últimos esquemas separan el entrenamiento de la validación para ofrecer estimaciones más fiables y constituyen una referencia metodológica más sólida para futuras pruebas, aunque continúan procesando la totalidad del fotograma.

La velocidad de procesamiento en el seguimiento peatonal se ve afectada tanto por la complejidad del algoritmo de rastreo como por la resolución de entrada, según se resume en la Tabla 1. En las pruebas de Li y Huang \cite{li2024} sobre el conjunto MOT17 manteniendo fijo el detector YOLOX, el rastreador ligero ByteTrack sostuvo una alta tasa de cuadros por segundo, mientras que incorporar extracción de apariencia con StrongSORT, BoT-SORT o siluetas de marcha incrementó ligeramente la precisión de seguimiento, pero redujo de forma pronunciada la velocidad. Por su parte, Zhang et al. \cite{zhang2023} evaluaron el sistema JYLD sobre el conjunto MOT20 y registraron una caída cercana a la mitad en los cuadros por segundo al elevar la resolución de entrada con el mismo rastreador. Como estas pruebas cambian de rastreador o aumentan la resolución, la disminución de FPS se debe a esos factores y no puede atribuirse solo a procesar zonas periféricas, por lo que aún se requiere medir por separado el efecto de acotar el área vigilada.

\begin{table}[!ht]
\centering
\scriptsize
\setlength{\tabcolsep}{5pt}
\renewcommand{\arraystretch}{1.1}
\begin{tabular}{|l|c|c|l|c|c|c|}
\hline
{\bfseries Estudio} & {\bfseries Conjunto} & {\bfseries Detector} & {\bfseries Rastreador} & {\bfseries Resolución} & {\bfseries MOTA} & {\bfseries Velocidad}\\
\hline
Li y Huang \cite{li2024} & MOT17 & YOLOX & ByteTrack & 1088 $\times$ 608 px & 76.7\,\% & 29.6 FPS\\
\hline
Li y Huang \cite{li2024} & MOT17 & YOLOX & StrongSORT & 1088 $\times$ 608 px & 77.1\,\% & 7.5 FPS\\
\hline
Li y Huang \cite{li2024} & MOT17 & YOLOX & BoT-SORT & 1088 $\times$ 608 px & 78.5\,\% & 4.5 FPS\\
\hline
Li y Huang \cite{li2024} & MOT17 & YOLOX & Siluetas de marcha & 1088 $\times$ 608 px & 78.0\,\% & 10.1 FPS\\
\hline
Zhang et al. \cite{zhang2023} & MOT20 & YOLOA-S & JYTracker & 512 $\times$ 928 px & 86.7\,\% & 20.5 FPS\\
\hline
Zhang et al. \cite{zhang2023} & MOT20 & YOLOA-S & JYTracker & 800 $\times$ 1440 px & 89.3\,\% & 10.8 FPS\\
\hline
\end{tabular}
\vspace{2.5mm}
\caption{Características de rendimiento y velocidad de procesamiento según el rastreador y la resolución de entrada.}\label{tab1}
\end{table}



La restricción espacial del área vigilada constituye otra estrategia para acotar el seguimiento peatonal a zonas de interés específicas, cuyas propuestas se resumen en la Tabla 2. En sistemas basados en cámaras convencionales, la estrategia predominante consiste en definir líneas de cruce o polígonos virtuales fijos \cite{fedosov2021,vu2026}. Sin embargo, estos filtros estáticos son vulnerables a las condiciones cambiantes del entorno, omitiendo eventos ante cruces simultáneos, sombras y oclusiones \cite{vu2026}, limitación que otros autores mitigan mediante la calibración geométrica de la región de interés para estimar distancias y detectar peatones lejanos con alta precisión \cite{lee2026}. Por otra parte, la delimitación espacial puede trasladarse al dominio tridimensional mediante sensores LiDAR para aislar zonas de riesgo a partir de nubes de puntos \cite{zhao2025}, prescindiendo de las variaciones de iluminación aunque con un costo de implementación restrictivo. Pese a estas alternativas, los métodos ópticos verifican las zonas delimitadas recién después de analizar el fotograma completo \cite{lee2026,fedosov2021,vu2026}, ya sea mediante estabilización previa en tomas aéreas \cite{jimenez2026}, máscaras estáticas sobre la vía \cite{yang2024} o estimación de pose sin recorte efectivo \cite{thosar2025}, mientras que las soluciones tridimensionales demandan equipamiento especializado ajeno a las cámaras estándar \cite{zhao2025}, dejando pendiente reducir la superficie evaluada por el detector durante la videovigilancia.

\begin{table}[!ht]
\centering
\scriptsize
\setlength{\tabcolsep}{5pt}
\renewcommand{\arraystretch}{1.1}
\begin{tabular}{|p{2.2cm}|p{2.3cm}|p{3.0cm}|c|c|}
\hline
{\bfseries Estudio} & {\bfseries Arquitectura} & {\bfseries Delimitación espacial} & {\bfseries Precisión o exactitud} & {\bfseries FPS}\\
\hline
Fedosov et al. \cite{fedosov2021} & YOLOv3 & Cruce peatonal manual fijo & 31--33\,\% AP & No reporta\\
\hline
Vu et al. \cite{vu2026} & YOLOv11, DeepSORT & Polígonos virtuales fijos & 66.0\,\% mAP & 10.6 FPS\\
\hline
Lee y Park \cite{lee2026} & YOLOv12, DeepSORT & Región de interés calibrada & 96.6\,\% & 17.2 FPS\\
\hline
Zhao et al. \cite{zhao2025} & EFT-RCNN & Área de peligro punto-región & 92.9\,\% & 28.56 FPS\\
\hline
\end{tabular}
\vspace{2.5mm}
\caption{Comparación de enfoques basados en regiones de interés y delimitación espacial.}\label{tab2}
\end{table}

A partir de la evidencia expuesta y de los resultados sintetizados en las Tablas 1 y 2, sostener una alta precisión en la detección del merodeo sin degradar la tasa de cuadros por segundo ni elevar el tiempo de respuesta del sistema requiere articular componentes específicos en las etapas de detección y rastreo peatonal. En esta tarea, resulta conveniente emplear detectores de una sola etapa de la familia YOLO junto con algoritmos de rastreo livianos como ByteTrack o DeepSORT \cite{lee2026,li2024,vu2026,zhang2023}, evitando extractores pesados de apariencia en toda la escena que reducen de forma drástica los cuadros por segundo \cite{li2024}. Para superar la limitación de las zonas fijas que solo descartan detecciones después de examinar el fotograma completo durante el monitoreo \cite{lee2026,fedosov2021,vu2026}, se plantea evaluar cercos virtuales dinámicos que permitan al usuario definir y reajustar manualmente los límites del área de interés según los requerimientos de vigilancia de la escena, concentrando el análisis exclusivamente en dicha región activa. Finalmente, dentro del área delimitada se pueden integrar criterios de tiempo de permanencia y cambios de dirección para identificar el merodeo \cite{jiang2022,li2024,m2022,zhang2023}, tomando como referencia los esquemas de validación de Waddenkery y Soma \cite{waddenkery2024} y Mehmood \cite{mehmood2021} para la exactitud, junto con los registros de Li y Huang \cite{li2024}, Zhang et al. \cite{zhang2023} y Lee y Park \cite{lee2026} sobre conjuntos de datos como UCF-Crime, IITB-Corridor o MOT20 para evaluar el costo computacional a partir de la tasa de cuadros por segundo y la latencia de procesamiento.


%
% ---- Bibliography ----
%
% BibTeX users should specify bibliography style 'splncs04'.
% References will then be sorted and formatted in the correct style.
%
% \bibliographystyle{splncs04}
% \bibliography{mybibliography}
%
\begin{thebibliography}{20}

\bibitem{choudhry2023}
Choudhry, N., Abawajy, J., Huda, S., Rao, I.: A Comprehensive Survey of Machine Learning Methods for Surveillance Videos Anomaly Detection. IEEE Access \textbf{11}, 114680--114713 (2023). \doi{10.1109/ACCESS.2023.3321800}

\bibitem{lee2026}
Lee, Y.H., Park, M.: A Real-Time Pedestrian Situation Detection Method Using CNN and DeepSORT with Rule-Based Analysis for Autonomous Mobility. Electronics \textbf{15}(3), 532 (2026). \doi{10.3390/electronics15030532}

\bibitem{mehmood2021}
Mehmood, A.: Abnormal Behavior Detection in Uncrowded Videos with Two-Stream 3D Convolutional Neural Networks. Applied Sciences \textbf{11}(8), 3523 (2021). \doi{10.3390/app11083523}

\bibitem{nayak2019}
Nayak, R., Behera, M.M., Girish, V., Pati, U.C., Das, S.K.: Deep Learning Based Loitering Detection System Using Multi-Camera Video Surveillance Network. In: 2019 IEEE International Symposium on Smart Electronic Systems (iSES) (Formerly iNiS), pp. 215--220 (2019). \doi{10.1109/iSES47678.2019.00055}

\bibitem{vinod2024}
Vinod, J., T M, D., Rajesh, R.: Interactive YOLO-based Object Detection using a Polygonal Region of Interest for Airborne Surveillance Applications. In: 2024 IEEE Space, Aerospace and Defence Conference (SPACE), pp. 901--905 (2024). \doi{10.1109/SPACE63117.2024.10668012}

\bibitem{wahyono2023}
Wahyono, Harjoko, A., Dharmawan, A., Adhinata, F.D., Kosala, G., Jo, K.-H.: Loitering Detection Using Spatial-Temporal Information for Intelligent Surveillance Systems on a Vision Sensor. Journal of Sensor and Actuator Networks \textbf{12}(1), 9 (2023). \doi{10.3390/jsan12010009}

\bibitem{zhao2025}
Zhao, A., Zhao, Y., Zheng, Q.: Robust Pedestrian Detection and Intrusion Judgment in Coal Yard Hazard Areas via 3D LiDAR-Based Deep Learning. Sensors \textbf{25}(18), 5908 (2025). \doi{10.3390/s25185908}

\bibitem{fedosov2021}
Fedosov, V.P., Ibadov, S.R., Ibadov, R.R., Kucheryavenko, S.V.: Method For Detecting Violation at a Pedestrian Crossing Using a Convolutional Neural Network. In: 2021 Radiation and Scattering of Electromagnetic Waves (RSEMW), pp. 451--454 (2021). \doi{10.1109/RSEMW52378.2021.9494089}

\bibitem{jeoung2026}
Jeoung, J., Hong, T., Park, J., Jeong, K.: Automatic discovery of semantic ROIs for vision-based occupancy analysis in corridors. Building and Environment \textbf{298}, 114656 (2026). \doi{10.1016/j.buildenv.2026.114656}

\bibitem{jiang2022}
Jiang, Z., Song, G., Qian, Y., Wang, Y.: A deep learning framework for detecting and localizing abnormal pedestrian behaviors at grade crossings. Neural Computing and Applications \textbf{34}(24), 22099--22113 (2022). \doi{10.1007/s00521-022-07660-0}

\bibitem{jimenez2026}
Jiménez-Frías, G.A., Hernandez-Lopez, F.J., Batard, T., Múniz-Sánchez, V., Hayet, J.-B., Forti-Sosa, S., Pérez-Pech, E.: Improving the Automatic Counting in a Region of Interest in Videos of Urban Traffic Taken by Drones Using Video Stabilization. Computación y Sistemas \textbf{30}(1) (2026). \doi{10.13053/CyS-30-1-6328}

\bibitem{li2024}
Li, H., Huang, X.: Spatiotemporal heterogeneous information fusion model for loitering anomaly detection. Signal, Image and Video Processing \textbf{18}(8--9), 5733--5742 (2024). \doi{10.1007/s11760-024-03267-1}

\bibitem{m2022}
M, S., S, G., Bharathy, S.: Loitering Detection in Home Surveillance System. In: 2022 10th International Conference on Emerging Trends in Engineering and Technology - Signal and Information Processing (ICETET-SIP-22), pp. 1--6 (2022). \doi{10.1109/ICETET-SIP-2254415.2022.9791651}

\bibitem{thosar2025}
Thosar, D., Thosar, R., Sharma, L., Chanore, M., Belkhede, P., Pohanerkar, S.: SmartSurveil: AI-based Intrusion and Suspicious Behavior Detection in CCTV. In: 2025 3rd DMIHER International Conference on Artificial Intelligence in Healthcare, Education and Industry (IDICAIHEI), pp. 1--6 (2025). \doi{10.1109/IDICAIHEI65991.2025.11379088}

\bibitem{vu2026}
Vu, T.C., Nguyen, D.T., Dinh, L.Q., Nguyen, M.D., Setiadi, D.R.I.M., Nguyen, M.T.: Deep-Learning-Based Multi-Camera Framework for Indoor Human Detection and Presence Management. Future Internet \textbf{18}(8), 435 (2026). \doi{10.3390/fi18080435}

\bibitem{waddenkery2024}
Waddenkery, N., Soma, S.: Loitering Based Human Crime Detection in Video Surveillance using Beluga Whale Adam Dingo Optimizer and Deep Convolutional Neural Network. The International Arab Journal of Information Technology \textbf{21}(3) (2024). \doi{10.34028/iajit/21/3/1}

\bibitem{yang2024}
Yang, X., Li, J.Q., Zhan, Y., Yu, W.: Real-time automated deep learning based railroad trespassing violation detection and tracking at highway-rail grade crossing. Intelligent Transportation Infrastructure \textbf{3}, liae003 (2024). \doi{10.1093/iti/liae003}

\bibitem{zhang2023}
Zhang, Y., Wang, J., Wang, X., Wang, Y., Yu, S.: Research and Application of Loitering Detection Based on Deep Learning. In: 2023 10th International Conference on Dependable Systems and Their Applications (DSA), pp. 339--346 (2023). \doi{10.1109/DSA59317.2023.00048}

\bibitem{cho2025}
Cho, J., Kang, Y.: Context-Aware Anomaly Detection of Pedestrian Trajectories in Urban Back Streets Using a Variational Autoencoder. ISPRS International Journal of Geo-Information \textbf{14}(11), 438 (2025). \doi{10.3390/ijgi14110438}

\bibitem{page2021}
Page, M.J., McKenzie, J.E., Bossuyt, P.M., Boutron, I., Hoffmann, T.C., Mulrow, C.D., et al.: The PRISMA 2020 statement: an updated guideline for reporting systematic reviews. BMJ \textbf{372}, n71 (2021). \doi{10.1136/bmj.n71}

\end{thebibliography}
\end{document}